% s-243-p243-m003; original PDF 243, 244
\Aphorism{s-243-p243-m003}{338}
Schliesslich ist doch die Natur, die ganz simple Licht-, Luft-, Wald- und Feldnatur die einzige Kunstmeisterin — oder nicht? Unsere vielgepriesene Seelenschilderei wäre zum Davonlaufen langweilig, wenn sie nicht Linien und Farben von aussen her entlehnte und alles Innerliche ins Morphologische, Decorative, Landschaftliche, überhaupt Bildliche und Malerische übersetzte. Die Psychophysik bedarf der Körperphysik, um geniessbar zu sein, und die bisher so arg spiritualisirte „Seele“ des Menschen, die unverkennbar zum unerlaubten Genre gehörte, wird immer interessanter werden, je enger und hingebender sie sich an seine Sinnlichkeit anschliesst. Tüchtig darauf los gelebt, geformt, genervt und gemuskelt! und weniger gedacht, begriffelt, gefrömmelt und gesittet! Seien wir Hylozoisten! keine Zahlen- und Geistesanbeter! Der Geist ist nur eine Species der Natur, und keine besonders wohlgerathene. Wenn ich nach meinem Geschmack die Menschheit zu verbessern hätte, ich schenkte ihr lieber eine neue physiologische Farbe, etwas Infrarothes oder Ultraviolettes, als hundert theoretische Abhandlungen über das Sehn und die Farben, — lieber eine neue Synthesis als die analytische Entzifferung von hundert alten Synthesen. Wieviel ist der objectiven Natur noch abzuringen, die uns von ihren zahllosen Schwingungen und Oscillationen den grössten Theil vorenthält! Für ganze Breiten undulatorischer Vorgänge fehlt uns ja immer noch die specifische Sinnesqualität! Das reiche Gebiet zwischen den höchsten akustischen und tiefsten thermischen Tönen ist noch zu erobern, nicht minder alles, was über achthundert Billionen Schwingungen in der Secunde macht. Das offene Meer der „Realität“ liegt vor uns; auf! ihr Entdecker und Corsaren!
\par

